% ---------
%  Compile with "pdflatex hw0".
% --------
%!TEX TS-program = pdflatex
%!TEX encoding = UTF-8 Unicode
%!TEX spellcheck = English (Aspell)

\documentclass[11pt]{article}
\usepackage{jeffe,handout,graphicx}
\usepackage[utf8]{inputenc}		% Allow some non-ASCII Unicode in source

\def\Cdot{\mathbin{\text{\normalfont \textbullet}}}
\def\Sym#1{\texttt{\color{BrickRed}#1}}
\def\BlueSym#1{\textbf{\texttt{\color{RoyalBlue}#1}}}

% =========================================================
\begin{document}

% ---------------------------------------------------------


\noindent Given a string $w$ and a symbol $a$, let \emph{delete}$(a, w)$ be the string $w$ with all instances of $a$ removed. For example, \emph{delete}(\Sym{z}, \Sym{jazzy}) is \Sym{jay} and \emph{delete}(\Sym{1},\Sym{00101110}) is \Sym{0000}.

\begin{enumerate}[(a)]
	\item 

Give a recursive function to compute \emph{delete}.


\begin{Solution}
Let $a$ be a symbol and $w$ a string.
(Hint: there are at least two cases based on the recursive definition of strings, but you may need more than two!)
\[
	\emph{delete}(a,w) := \begin{cases}
		 & \text{if }\\[0.25ex]
		 & \text{if }\\
	\end{cases}
\]
\end{Solution}

\begin{Sources}~ % note the tilde to add a line break
\begin{itemize}
    \item  Bob Smith helped me. He took this course at UM in Fall 2020. He suggested that I needed two separate recursive cases.
    \item I looked at the slides at \url{https://chengjianglong.com/slides/DS_Lecture_22.pdf} for another perspective on recursive structures.
	\item I used ChatGPT. Here is the full transcript of my conversation:\\
	 \url{https://chatgpt.com/share/6970e55f-e594-800e-88e3-c5eaa8caf7ce}

\end{itemize}
\end{Sources}

\newpage

\item 
Prove that $\emph{delete}(a, x \Cdot y) = \emph{delete}(a, x) \Cdot \emph{delete}(a,y)$ for all strings $x,y$ and all symbols $a$.

\begin{Solution}
\parindent0pt
Let $x,y$ be arbitrary strings and $a$ an arbitrary symbol.

Assume... (what's the inductive hypothesis?)

There are ?? cases to consider:
\begin{itemize}
\item
If $w = \e$, then
\begin{align*}
	\emph{delete}(a,x \Cdot y)
	&= \emph{delete}(a,\e \Cdot y)	& \text{because ?} \\
	&\\
	&\\
	&\\
	&\\
	&= \emph{delete}(a, \e) \Cdot \emph{delete}(a, y)	& \text{because ?} \\
	&= \emph{delete}(a, x) \Cdot \emph{delete}(a, y)	& \text{because $w = \e$} \\
\end{align*}
\item
...you fill in the rest of the cases
\end{itemize}
In all cases, we conclude that $\emph{delete}(a, x \Cdot y) = \emph{delete}(a, x) \Cdot \emph{delete}(a,y)$.
\end{Solution}
\unskip

\begin{Sources}~ % note the tilde to add a line break
\begin{itemize}
	\item Same ChatGPT conversation as above. Here is the full transcript of my conversation:\\
	 \url{https://chatgpt.com/share/6970e55f-e594-800e-88e3-c5eaa8caf7ce}

\end{itemize}
\end{Sources}

\newpage
\item Recall the function $\#(a, w)$ from problem session 1, which returns the number of occurrences of the symbol $a$ in string $w$. For any string $w \in \{0,1\}^*$, prove that $\abs{\emph{delete}(1,w)} = \#(0, w)$.

\begin{Solution}
\parindent0pt
Let $w \in \{0,1\}^* $.

\vspace{6in}

In all cases, we conclude that  $\abs{\emph{delete}(1,w)} = \#(0, w)$.
\end{Solution}

\begin{Sources}
None.
\end{Sources}


\end{enumerate}



\end{document}
